%!TEX root = hems_proj.tex
%%%%%%%%%%%%%%%%%%%%%%%%%%%%%%%%%%%%%%%%%%%%%%%%%%%%%%%%%%%%%%%%%%%%%%%
\chapter{Alkalmazott Algoritmusok}\label{ch:ELM}
%%%%%%%%%%%%%%%%%%%%%%%%%%%%%%%%%%%%%%%%%%%%%%%%%%%%%%%%%%%%%%%%%%%%%%%

\begin{osszefoglal}
Négy metaheurisztikát és egy hibrid eljárást alkalmaztunk: GA, PSO, GWO és GA-PSO hibrid. A fejezet bemutatja mindegyik algoritmus működési elvét, paraméterezését és a HEMS-feladatra való alkalmazhatóságát.
\end{osszefoglal}

\section{Genetikus Algoritmus (GA)}

A genetikus algoritmusok alapötletét Holland klasszikus munkája alapozta meg \cite{HollandGA1975}. A módszer a természetes szelekció elvét alkalmazza: egy populációnyi megoldásjelöltet tart fenn, amelyek generációról generációra fejlődnek keresztezés és mutáció révén. Fő előnye, hogy jól kezeli a diszkrét változókat és robusztus megoldást nyújt változékony keresési térben \cite{Ismail2014}, bár hátránya a PSO-hoz képest lassabb konvergencia.

Az implementált algoritmus működése a következő. Elsőként $N$ darab véletlen egyedet generálunk, amelyek a kezdeti populációt alkotják, majd minden egyedhez kiszámítjuk a fitness értéket a célfüggvény alapján. A fő ciklus $G$ generáción keresztül fut. Minden iterációban az elitizmus elve alapján a legjobb két egyed automatikusan továbbjut. A szülőket tournament módszerrel választjuk ki, majd aritmetikai átlagolással és zaj hozzáadásával keresztezzük őket. A sokszínűség fenntartása érdekében Gauss-eloszlású véletlen módosítással mutációt hajtunk végre, majd az algoritmus visszatér a legjobb megoldással. A mérések alapján az algoritmus átlagosan 15--20\%-os költségcsökkenést eredményez az optimalizálás nélküli alapvonalhoz képest.

\section{Particle Swarm Optimization (PSO)}

A PSO eredeti formáját Kennedy és Eberhart vezette be 1995-ben, madárrajok és haliskolák kollektív viselkedésének modellezéséből kiindulva \cite{KennedyEberhartPSO1995}. Az algoritmus alapegysége a részecske (particle), amely egy lehetséges megoldást képvisel a keresési térben. Minden részecskének van egy aktuális pozíciója $\vec{x}_i$ és egy sebessége $\vec{v}_i$, amelyek iterációról iterációra frissülnek.

A sebesség-frissítési egyenlet három tagból áll. Az első a tehetetlenségi tag, amely a részecske korábbi mozgásirányát tartja fenn: $w \cdot \vec{v}_i^k$. A második a kognitív tag, amely a részecske saját eddigi legjobb pozíciója ($\vec{p}_{best,i}$) felé húzza: $c_1 r_1 (\vec{p}_{best,i} - \vec{x}_i^k)$. A harmadik a szociális tag, amely a teljes raj által talált legjobb pozíció ($\vec{g}_{best}$) felé irányít: $c_2 r_2 (\vec{g}_{best} - \vec{x}_i^k)$. Összefoglalva:
\begin{equation}
\vec{v}_i^{k+1} = w \cdot \vec{v}_i^k + c_1 r_1 (\vec{p}_{best,i} - \vec{x}_i^k) + c_2 r_2 (\vec{g}_{best} - \vec{x}_i^k)
\label{eq:pso_velocity}
\end{equation}
ahol $r_1, r_2 \in [0,1]$ egyenletesen eloszló véletlen számok. A pozíció frissítése ezután:
\begin{equation}
\vec{x}_i^{k+1} = \vec{x}_i^k + \vec{v}_i^{k+1}
\label{eq:pso_position}
\end{equation}

A tehetetlenségi súly $w$ szabályozza a globális és lokális keresés egyensúlyát: nagy $w$ értéknél a részecskék szélesebb területet pásztáznak (exploráció), kis $w$ értéknél pedig a már megtalált jó megoldások közelében finomhangolnak (exploitáció). A $c_1$ és $c_2$ kognitív és szociális gyorsulási konstansok azt szabályozzák, hogy a részecske mennyire bízik saját tapasztalatában versus a raj kollektív tudásában. A PSO előnye a gyors konvergencia és az egyszerű implementáció, hátránya azonban, hogy zajos keresési térben hajlamos hamis lokális optimumba ragadni, mivel a $\vec{g}_{best}$ komponens gyorsan magához vonzza az összes részecskét, csökkentve a populáció sokszínűségét.

\section{Grey Wolf Optimizer (GWO)}

A GWO algoritmust Mirjalili és szerzőtársai javasolták 2014-ben, a szürke farkasok hierarchikus csoportvadászatát modellezve \cite{MirjaliliGWO2014}. Az algoritmus négy rangszintet különböztet meg: az $\alpha$ farkas a legjobb megoldást képviseli, a $\beta$ a másodikat, a $\delta$ a harmadikat, az $\omega$ farkasok pedig a populáció többi tagját alkotják. Ez a hierarchia a keresés irányítási mechanizmusaként működik.

Minden iterációban az $\alpha$, $\beta$ és $\delta$ farkasok pozíciója alapján frissül az összes többi egyed helyzete. Először meghatározzuk az egyes vezérfarkasokhoz való távolságot:
\begin{align}
\vec{D}_\alpha &= |\vec{C}_1 \cdot \vec{X}_\alpha - \vec{X}| \\
\vec{D}_\beta  &= |\vec{C}_2 \cdot \vec{X}_\beta  - \vec{X}| \\
\vec{D}_\delta &= |\vec{C}_3 \cdot \vec{X}_\delta - \vec{X}|
\end{align}
Ezután kiszámítjuk az egyes vezérek által javasolt új pozíciókat:
\begin{align}
\vec{X}_1 &= \vec{X}_\alpha - \vec{A}_1 \cdot \vec{D}_\alpha \\
\vec{X}_2 &= \vec{X}_\beta  - \vec{A}_2 \cdot \vec{D}_\beta  \\
\vec{X}_3 &= \vec{X}_\delta - \vec{A}_3 \cdot \vec{D}_\delta
\end{align}
Az egyед végső új pozícióját a három javaslat átlagaként kapjuk:
\begin{equation}
\vec{X}(t+1) = \frac{\vec{X}_1 + \vec{X}_2 + \vec{X}_3}{3}
\end{equation}

Az $\vec{A}$ és $\vec{C}$ vektorok az iteráció előrehaladásával változnak: az $a$ paraméter lineárisan csökken 2-ről 0-ra, ami azt eredményezi, hogy a keresés kezdetben széles feltárást (exploráció), majd fokozatosan finomhangolást (exploitáció) végez. A GWO előnye, hogy kevés hangolható paramétert igényel és jól egyensúlyoz az exploráció és exploitáció között, különösen szűkebb, korlátosabb keresési terekben.

\section{Hibrid GA-PSO}

A GA és PSO kombinálása ismert hibridizációs irány az energiarendszeres optimalizálásban \cite{KumarGA}. A hibrid algoritmus két egymást követő fázisban dolgozik. Az első fázisban -- az iterációk első 60\%-ában -- a PSO végzi a gyors feltérképezést a sebesség-frissítési egyenletek segítségével, kiaknázva a raj kollektív tudását. A második fázisban a GA veszi át az irányítást: a PSO által megtalált legjobb megoldás köré épített populációból indul ki, majd keresztezéssel és mutációval biztosítja a megoldások finomhangolását és a diverzitás fenntartását. A két fázis közötti váltásnál az úgynevezett seed injection technikát alkalmazzuk: a PSO legjobb megoldása bekerül a GA kiindulási populációjába, a populáció fele pedig annak véletlen perturbációiból áll, így a GA jó helyről indul, de elegendő variációja marad a további fejlődéshez.

\section{Komplexitás és paraméterérzékenység}

Valamennyi vizsgált algoritmus populációalapú keresést végez, ezért a végrehajtási idő első közelítésben a populációméret ($N$), az iterációszám ($G$) és a döntési dimenzió ($d$) szorzatával arányos. A projekt jelenlegi problémájában $d=26$: 24 akkumulátorvezérlési változó és 2 eszközindítási idő szerepel a döntési vektorban:
\begin{equation}
\mathcal{O}(N \cdot G \cdot d)
\end{equation}
A paraméterérzékenységi vizsgálat alapján a nagyobb populáció stabilabb keresést ad, de megnöveli a futási időt. A PSO alacsony futási ideje mellett hajlamosabb korai stagnálásra zajos bemeneteken. A GWO kevés paraméterrel is erős kiindulópontot ad, különösen korlátozott keresési térben. A hibrid modell csak akkor használja ki valóban az előnyeit, ha a PSO-fázis nem túl rövid és a GA-fázis elegendő finomhangolási teret kap.

\section{A metaheurisztikák szerepe a terheléseltolásban}

A terheléseltolásos HEMS-feladatokban a keresés nem csak költségminimalizálásról szól, hanem a megvalósíthatóság és a működési komfort közti egyensúlyról is \cite{Fadaee2012}. A GA képes több, eltérő kezdetű megoldást fenntartani; a PSO gyorsan reagál a jobb egyedekre; a GWO kiegyensúlyozott, hierarchikus keresést biztosít; a hibrid pedig egyesíti a feltáró és finomító fázisokat. Az eredmények alapján különösen fontos tanulság, hogy a megoldás minősége nem lineárisan javul a futási idő növelésével -- egyes esetekben a gyors algoritmusok is elérnek jó lokális megoldásokat, de a nagyobb változékonyságú napokon a robusztusabb keresési stratégia válik előnyössé.